\def\topic{Supremum \& Infimum}
\ifdefined\MyHeaderIncluded
\else
\def\qwertyflag{true}

\ifdefined\MyHeaderIncluded

\else
  \newcommand{\foobar}{FooBar}
  \def\MyHeaderIncluded{true}
% if is till the end of the file

\documentclass[10pt,a4paper]{article}
\usepackage[utf8]{inputenc}
%\usepackage[german]{babel}
%\usepackage[cyrillic]{babel}
\usepackage[T1]{fontenc}
%\usepackage[T2A]{fontenc}
%\usepackage[english, russian]{babel}
\usepackage{amsmath}
\usepackage{amsfonts}
\usepackage{amssymb}
\usepackage{amsthm}
\usepackage{mathtools} %coloneqq

\usepackage{bussproofs} %

%\renewcommand{\thesubsection}{\arabic{subsection}}
%\newtheorem{thm}{Theorem}[section]
%\renewcommand{\thethm}{\arabic{subsection}.\arabic{thm}}

\usepackage{thmtools}
\usepackage{hyperref}
\hypersetup{
    colorlinks=true,
    linkcolor=blue,
    filecolor=magenta,      
    urlcolor=cyan,
    pdftitle={Overleaf Example},
    %pdfpagemode=FullScreen,
    }
%\usepackage{stix-type1}%\turnediota
\theoremstyle{definition}

\newtheorem{constructor}{Constructor}[section]
\newtheorem{iaxiom}{Axiom}
\newtheorem{axiom}{Axiom}
\newtheorem*{axiom*}{Axiom}
\newtheorem{definition}{Definition}[section]
\newtheorem*{definition*}{Definition}
\newtheorem{remark}{Remark}[section]
\newtheorem*{remark*}{Remark}
\newtheorem{recall}{Recall}[section]
\newtheorem*{recall*}{Recall}

\newtheorem{lemma}{Lemma}[section]
\newtheorem{theorem}{Theorem}[section]
\newtheorem*{theorem*}{Theorem}
\newtheorem{note}{Note}[section]
\usepackage{etoolbox}

%\newcommand*{\@Dashv}[2]{%
%  \reflectbox{$\m@th#1#2$}%
%}
\newcommand*{\separation}[3]{% 1)binder, 2)set, 3)condition
  \{#1\in #2: #3\}
}

% https://www.freecodecamp.org/news/simple-regex-tricks-for-beginners-3acb3fa257cb/
%  \\mathcal\{T\} |---> \topology
%  \\mathcal\{([^\}]*)\}   |--->  $1
\newcommand{\AND}{\ \ \&\ \ }
\newcommand{\any}{\mathrm{\ any}}
\newcommand{\borel}{\mathcal{B}}
\newcommand{\can}{\mathrm{can}}
\newcommand{\canonize}{\mathrm{canonize}}
\newcommand{\closure}{\mbox{cl}}
\newcommand{\complex}{\mathbb{C}}
\newcommand{\converges}{\longrightarrow}
\newcommand{\countable}{\mbox{\ countable}}
\newcommand{\ldefined}{\mbox{\ defined}}
\newcommand{\dom}{\mbox{dom}}
\newcommand{\Dom}{\mbox{Dom}}
\newcommand{\Done}{\noindent \textbf{Done.}\\}
\newcommand{\lexists}{\mbox{\ exists}}
\newcommand{\lforall}{\mbox{\ forall}}
\newcommand{\lunique}{\mbox{\ unique}}
\newcommand{\family}{\mathrm{\ family}}
\newcommand{\Fm}{\mathrm{Fm}}
\newcommand{\FuncSymb}{\mathrm{FuncSymb}}
\newcommand{\function}{\mathrm{\ function}}
\newcommand{\functional}{\mbox{\ functional}}
\newcommand{\IFF}{\ \Longleftrightarrow\ }
\renewcommand{\iff}{\ \longleftrightarrow\ }
\newcommand{\IMPLIES}{\ \Longrightarrow\ }
\renewcommand{\implies}{\ \longrightarrow\ }
\newcommand{\Inf}{\mathrm{Inf}}
\newcommand{\injective}{\mbox{\ injective}}
\newcommand{\injection}{\mbox{\ injection}}
\newcommand{\interior}{\mbox{int}}
\newcommand{\intersection}{\bigcap}
\newcommand{\isempty}{\mathrm{\ empty}}
\newcommand{\LB}{\mathrm{LB}}
\newcommand{\nghbset}{\mathcal{N}}
\newcommand{\nonempty}{\mathrm{\ nonempty}}
\newcommand{\Ord}{\mathrm{Ord}}
\newcommand{\power}{\mathcal{P}}
\newcommand{\pr}{\mathrm{pr}}
\newcommand{\fst}{\mathrm{fst}}
\newcommand{\snd}{\mathrm{snd}}
\newcommand{\PredSymb}{\mathrm{PredSymb}}
\newcommand{\Proof}{\noindent \textbf{Proof.}\\}
\newcommand{\ran}{\mathrm{ran}}
\newcommand{\Ran}{\mathrm{Ran}}
\newcommand{\real}{\mathbb{R}}
\newcommand{\relation}{\mathrm{\ relation}}
\newcommand{\set}{\mathrm{\ set}}
\newcommand{\singleton}{\mathrm{\ singleton}}
\newcommand{\semantics}{\mathrm{semantics}}
\newcommand{\semialgebra}{\mathcal{S}}
\newcommand{\Sup}{\mathrm{Sup}}
\newcommand{\surjective}{\mbox{\ surjective\ on\ }}
\newcommand{\Term}{\mathrm{Term}}
\newcommand{\Type}{\mathrm{Type}}
\newcommand{\topology}{\mathcal{T}}
\newcommand{\total}{\mbox{\ total\ on\ }}
\newcommand{\transitive}{\mathrm{\ transitive}}
\newcommand{\union}{\bigcup}
\newcommand{\lTrue}{\mathrm{True}}
\newcommand{\lFalse}{\mathrm{False}}
\newcommand{\lUncertain}{\mathrm{Uncertain}}
\newcommand{\lif}{\mathrm{if}}
\newcommand{\lthen}{\mathrm{then}}
\newcommand{\lelse}{\mathrm{else}}
\newcommand{\lexception}{\mathrm{exception}}
\newcommand{\lundef}{\mathrm{undef}}
\newcommand{\univocal}{\mathrm{\ univocal}}
\newcommand{\univocalat}{\mathrm{\ univocal\ at\ }}
\newcommand{\UB}{\mathrm{UB}}
\newcommand{\Un}{\mathcal{U}}
\newcommand{\Var}{\mathrm{Var}}

\newcommand{\notvarphi}{\phi}
\renewcommand{\phi}{\varphi}

\newcommand{\comment}[1]{\quad\mbox{(#1)}}
% local commands
%\preto\section{\undef\blipo} %etoolbox
\preto\subsection{\undef\blipo} %etoolbox
% https://tex.stackexchange.com/questions/105684/limit-macro-newcommand-scope-to-section
\usepackage{graphicx}
\usepackage{authblk}

\setlength\parindent{0pt} % whole doc of \noindent
\usepackage{soul}
%\usepackage{mathgl}

\fi % guard
\ifdefined\topic
\else
\def\topic{No Titel}
\fi
\title{\topic}
\author[1]{Georgy Dunaev}
\begin{document}
\maketitle
\fi

%Some order theory.
\section{Suprema and infima}
The first construction retains the inherited order on the class of bounds.
This allows it to describe preorders, where a least element need not be
unique as an object, as well as partial orders, where it is unique when it
exists. For the image-based construction first, $(P,\leq)$ is a set-sized
ambient order and $S\subseteq P$. Its upper and lower sections are then
sets, so they can be collected by the displayed image notation. For a
proper-class carrier, use the predicate-based bound classes given later;
ordinary class images cannot collect proper classes as members. Whenever a carrier is restricted, the order relation is
restricted to that carrier as well.

\subsection{Up or down sets of an element of an order}%Upward/downward closed set}
\begin{definition}[up set of an element of an order]%(Upward closed set)
\[
\uparrow_{(P,\leq)}x\ceq \cc{y\in P}{x\leq y}
\]
\end{definition}
\begin{definition}[down set of an element of an order]%(Downward closed set)
\[
\downarrow_{(P,\leq)}x\ceq \cc{y\in P}{y\leq x}
\]
\end{definition}

\subsection{Ordered set of majorants or minorants of a subset of an order}
\begin{definition}[ordered set of majorants of a subset of an order]
\[
\mathrm{Ma}_{(P,\leq)}S\ceq \left(P\cap\intersection \uparrow_{(P,\leq)}[S],\leq\right)
\]
\end{definition}
\begin{definition}[ordered set of minorants of a subset of an order]
\[
\mathrm{Mi}_{(P,\leq)}S\ceq \left(P\cap\intersection \downarrow_{(P,\leq)}[S],\leq\right)
\]
\end{definition}
\subsection{Ordered set of greatest or least elements of an order}
\begin{definition}[ordered set of greatest elements of an order]
\[
\mathrm{Greatest}_{(P,\leq)}\ceq \mathrm{Ma}_{(P,\leq)}P
\]
\end{definition}
\begin{definition}[ordered set of least elements of an order]
\[
\mathrm{Least}_{(P,\leq)}\ceq \mathrm{Mi}_{(P,\leq)}P
\]
\end{definition}
\subsection{Ordered set of suprema or infima of a subset of an order}
\begin{definition}[ordered set of suprema of a subset of an order]
\[
\mathrm{Sup}_{(P,\leq)}S\ceq \mathrm{Least}_{\mathrm{Ma}_{(P,\leq)}S}
\]
\end{definition}
\begin{definition}[ordered set of infima of a subset of an order]
\[
\mathrm{Inf}_{(P,\leq)}S\ceq \mathrm{Greatest}_{\mathrm{Mi}_{(P,\leq)}S}
\]
\end{definition}
\subsection{Essential supremum}
Here the context changes to a measure space and a measurable real-valued
function $f$. This value depends only on the almost-everywhere equivalence
class of $f$: changing $f$ on a measurable null set does not change which
real numbers are essential upper bounds. Take the infimum in the extended real line, so that the
expression also covers an unbounded function: an empty set of real upper
bounds has infimum $+\infty$. The displayed condition means that $M$ is an
upper bound except on a set of measure zero.

\begin{definition}[essential supremum; it depends on a measure $\mu$]
\[
\esssup f
\ceq
\inf
\cc{M\in\real}{\mu(\cc{x\in \dom(f)}{f(x)>M})=0}
\]
%\mu\shortminus
%https://www.youtube.com/watch?v=7xKAFQJfyPY
\end{definition}

With this extended-real convention, a zero-measure carrier gives
$\esssup f=-\infty$: every real number is an essential upper bound.
The later $L^\infty$ norm instead takes the infimum of \emph{nonnegative}
bounds for $|f|$, so it is zero on a zero-measure space. In these
conventions its value is $\max\{0,\esssup|f|\}$; on a space of
positive measure this is simply $\esssup|f|$, since no negative number
can bound $|f|$ almost everywhere. This distinction concerns the range
in which the infimum is taken, not a failure of the norm axioms.

%\begin{theorem}
%\[
%(A\set)
%\implies
%\left(\union A\set\right)
%\]
%\end{theorem}
%\Proof
%(the proof needs theorems about ordinals done first)\\
%...\\
%\Done
\subsection{Some theorems of the order theory}

The following reformulation uses classes of bounds directly. An ordinary
relational image $R[S]$ describes points above \emph{some} member of $S$;
upper bounds must lie above \emph{every} member. This distinction motivates
the universal definitions below. They express the intended meaning of the
earlier image-based sketches.
\begin{definition}[upper and lower bounds]
\[
\UB_R(S)\ceq\cc{y\in P}{\forall x\in S,\ xRy},\qquad
\LB_R(S)\ceq\cc{y\in P}{\forall x\in S,\ yRx}.
\]
\end{definition}

Context and notation:\\
Here by $\leq$ we denote the standard relation "less or equal to" on real numbers.\\
$\leq$ for inequalities of numbers;\\
$\leqslant$ for inequalities of sets of numbers;\\
$R \ceq \leq$\\
$\leq = \geq^{-1}$\\
$\geq = \leq^{-1}$\\

\begin{remark}[why the middle class must be nonempty]
Define $A\leqslant B$ to mean that every element of $A$ is at most every element of $B$. If $A\leqslant B$ and $B\leqslant C$, transitivity of $\leq$ gives $A\leqslant C$ provided the middle class $B$ is nonempty. Without that hypothesis, both premises can hold vacuously.
\end{remark}
\Proof
Choose $b\in B$. For arbitrary $a\in A$ and $c\in C$, we have $a\leq b\leq c$, hence $a\leq c$. For the failure with an empty middle class, take $S<H$. The familiar wordplay illustrates the two vacuous premises:
\begin{quote}
Nothing is better than eternal happiness. \null\hfill ($\{H\}<\emptyclass$)\\
A sandwich is better than nothing. \null\hfill ($\emptyclass<\{S\}$)\\
Therefore, a sandwich is better than eternal happiness. \null\hfill ($\{H\}<\{S\}$)\\
\end{quote}
We obtained $H<S$ (by definition of an order on sets). The statement together with $S<H$ and transitivity on elements contradicts the irreflexivity of the order on elements.\\
\Done
\noindent

\begin{definition}[all-to-all comparison]
\[
(A \leqslant B)
\iffdef
(\forall a\in A \forall b\in B,\ a\leq b)
\]
\end{definition}
\begin{definition}[bound classes]
\[
\UB(S) \ceq \UB_{\leq}(S)
\]
\end{definition}
\begin{definition}[bound classes]
\[
\LB(S) \ceq \LB_{\leq}(S)
\]
\end{definition}
\begin{theorem}[comparison by bounds]
\[
(A\leqslant B)
\iff
(B\subseteq \UB(A))
\]
\end{theorem}
\begin{theorem}[comparison by bounds]
\[
(A\leqslant B)
\Longleftrightarrow
(A\subseteq \LB(B))
\]
\end{theorem}
\begin{definition}[bound classes]
\[
\Sup(S) = \UB(S)\cap \LB(\UB(S))
\]
\end{definition}
\begin{definition}[bound classes]
\[
\Inf(S) = \LB(S)\cap \UB(\LB(S))
\]
\end{definition}
\begin{definition}[descriptor for the unique extremum]
\[
\mbox{sup}(S) = (\iota x,\ x\in \mathrm{Sup}(S))
\]
\end{definition}
\begin{definition}[descriptor for the unique extremum]
\[
\mbox{inf}(S) \eqdef (\iota x,\ x\in \mathrm{Inf}(S))
\]
\end{definition}
A descriptor for a supremum or infimum denotes the desired extremum only
when the corresponding bound class has exactly one member. Its behavior
otherwise is governed by the descriptor convention chosen locally.
\begin{theorem}[comparison of existing extrema]
Assume the ambient order is a partial order, and both $\Sup(A)$ and
$\Inf(B)$ are nonempty.
\[
(A \leqslant B)
\Longrightarrow
(A\leqslant \Sup(A)\leqslant \Inf(B)\leqslant B)
\]
\end{theorem}
\begin{proof}
Let $s$ and $i$ be the unique elements of $\Sup(A)$ and $\Inf(B)$.
Every $b\in B$ is an upper bound of $A$, so the least-upper-bound
property gives $s\leq b$. Thus $s$ is a lower bound of $B$, and the
greatest-lower-bound property gives $s\leq i$. The outer comparisons
$a\leq s$ for $a\in A$ and $i\leq b$ for $b\in B$ follow from the
definitions of upper and lower bounds. This proves the displayed chain.
\end{proof}

\begin{theorem}[monotone convergence of a real sequence]
\[
(a:\mathbb{N}\rightarrow \real)
\land
(\{a_n\}_n \mbox{bdd above in }\real)
\land
((a_n)_n\mbox{ non--decreasing})
\implies
(a \converges \sup \{a_n\}_n)
\]
\end{theorem}
\begin{proof}
The range of the sequence is nonempty and bounded above, so completeness
of the real numbers supplies its real supremum $s$. For every $\varepsilon>0$,
$s-\varepsilon$ is not an upper bound; choose $N$ with
$s-\varepsilon<a_N$. Monotonicity gives
\[
s-\varepsilon<a_N\leq a_n\leq s\qquad(n\geq N).
\]
Hence $|a_n-s|<\varepsilon$ for $n\geq N$, which is exactly convergence
to $s$. This argument uses a real supremum, not an off-domain descriptor
value or the ordinal infimum convention from another section.
\end{proof}

\subsection{Bounds, closure, and fixed points}
We now return from the real-number examples to a set-sized partially
ordered set $P$. Bound classes are always restricted to this carrier.
The equivalence
\[
B\subseteq\UB(A)\quad\Longleftrightarrow\quad A\subseteq\LB(B)
\qquad(A,B\subseteq P)
\]
is the order instance of the relational polarity proved in the section
``The order structure behind glare.'' In that notation the two maps are
$p_P=\UB$ and $q_P=\LB$. In particular,
\[
c(A)=\LB(\UB(A))
\]
is extensive, monotone, and idempotent, and
$\UB(\LB(\UB(A)))=\UB(A)$. These statements include empty subsets:
$\UB(\emptyclass)=P$. When $s=\sup A$ exists, $c(A)=\downarrow s$.
Thus the bound construction already supplies an example of a closure
operator. We next establish a more general existence theorem, which
requires monotonicity but not extensivity or idempotence.

\begin{definition}[Complete lattice and fixed-point inequalities]
A set-sized partially ordered set $L$ is a complete lattice if every
subset $A\subseteq L$ has a supremum $\bigvee_L A$ and an infimum
$\bigwedge_L A$ in $L$, including the empty subset. Write
$\bot=\bigvee_L\emptyclass$ and $\top=\bigwedge_L\emptyclass$.
For a monotone function $f:L\to L$, put
\[
\begin{aligned}
\operatorname{Pre}(f)&=\{x\in L:f(x)\leq x\},\\
\operatorname{Post}(f)&=\{x\in L:x\leq f(x)\},\\
\operatorname{Fix}(f)&=\{x\in L:f(x)=x\}.
\end{aligned}
\]
The first two conditions are called pre-fixed and post-fixed,
respectively; this specifies the direction of the terminology.
\end{definition}

\begin{theorem}[Complete lattice of fixed points: Knaster--Tarski]
Let $f:L\to L$ be monotone on a set-sized complete lattice. Its fixed
points form a complete lattice in the inherited order. More explicitly,
for $A\subseteq\operatorname{Fix}(f)$ define
\[
\begin{aligned}
U_A&=\{x\in L:f(x)\leq x\text{ and }\forall a\in A,\ a\leq x\},\\
D_A&=\{x\in L:x\leq f(x)\text{ and }\forall a\in A,\ x\leq a\}.
\end{aligned}
\]
Then the join and meet in the fixed-point lattice are
\[
\bigvee_{\operatorname{Fix}(f)}A=\bigwedge_L U_A,
\qquad
\bigwedge_{\operatorname{Fix}(f)}A=\bigvee_L D_A.
\]
In particular, the least and greatest fixed points are
\[
\mu f=\bigwedge_L\operatorname{Pre}(f),\qquad
\nu f=\bigvee_L\operatorname{Post}(f).
\]
The element $\mu f$ lies below every pre-fixed point, and $\nu f$
lies above every post-fixed point.
\end{theorem}
\Proof
The set $U_A$ contains $\top$. Put $j=\bigwedge_L U_A$.
Each $a\in A$ is a lower bound of $U_A$, so $a\leq j$.
For each $x\in U_A$, monotonicity gives
$f(j)\leq f(x)\leq x$. Thus $f(j)\leq j$. It follows that
$f(f(j))\leq f(j)$, so $f(j)$ is pre-fixed. Also, for $a\in A$,
$a=f(a)\leq f(j)$. Hence $f(j)\in U_A$, and the infimum property
gives $j\leq f(j)$. Antisymmetry proves $f(j)=j$.
Finally, every fixed upper bound $b$ of $A$ belongs to $U_A$, so
$j\leq b$. Therefore $j$ is precisely the join among fixed points.

For completeness, the dual argument is as follows. The set $D_A$
contains $\bot$. Put $m=\bigvee_L D_A$. Every $a\in A$ is an upper
bound of $D_A$, so $m\leq a$. If $x\in D_A$, then
$x\leq f(x)\leq f(m)$, whence $m\leq f(m)$. Therefore
$f(m)\leq f(f(m))$, and for each $a\in A$ we have
$f(m)\leq f(a)=a$. Thus $f(m)\in D_A$ and $f(m)\leq m$.
Antisymmetry makes $m$ fixed. Every fixed lower bound $b$ of $A$
belongs to $D_A$, so $b\leq m$. This proves the meet assertion.

Taking $A=\emptyclass$ gives $U_A=\operatorname{Pre}(f)$ and
$D_A=\operatorname{Post}(f)$, yielding the last assertions as well as
nonemptiness of the fixed-point lattice. All bound sets used here are
subsets of $L$. The proof uses neither a transfinite iteration nor a
choice of approximating elements.
\Done

\begin{example}[A join can change when restricted to fixed points]
On the powerset lattice of $\{a,b,c\}$, with these three elements
distinct, define
\[
f(X)=\begin{cases}
X\cup\{c\},&\{a,b\}\subseteq X,\\
X,&\text{otherwise}.
\end{cases}
\]
This map is monotone: enlarging $X$ preserves its members and cannot
turn off an already satisfied condition $\{a,b\}\subseteq X$.
Both $\{a\}$ and $\{b\}$ are fixed, but their ambient join
$\{a,b\}$ is not fixed. Their join among fixed points is
$\{a,b,c\}$, since every fixed set containing both $a$ and $b$ must
also contain $c$. Thus the theorem does not assert that fixed points
are closed under the ambient joins and meets. Its displayed formulas
specify the bounds in their own inherited order.
\end{example}

\begin{corollary}[Least closed subset containing prescribed data]
Let $X$ be a set, $S\subseteq X$, and let
$\Phi:\mathcal P(X)\to\mathcal P(X)$ be monotone. There is a least
subset $C\subseteq X$ containing $S$ and satisfying $\Phi(C)\subseteq C$.
It is given by
\[
C=\bigcap\{Q\subseteq X:S\subseteq Q\text{ and }\Phi(Q)\subseteq Q\},
\qquad C=S\cup\Phi(C).
\]
\end{corollary}
\Proof
Apply the theorem in the inclusion-ordered complete lattice
$\mathcal P(X)$ to $f(Q)=S\cup\Phi(Q)$. Its pre-fixed points are
exactly the subsets in the displayed intersection, a nonempty family
because it contains $X$. The least fixed point is below every such
subset and is itself one of them. Its fixed-point equation is exactly
$C=S\cup\Phi(C)$.
\Done

This is the set-sized counterpart of the least-class closure theorem in
Section~\ref{sec:properties_preserved_intersections}. That theorem uses
MK comprehension over classes; it is not obtained by treating the
collection of all proper classes as a set-sized complete lattice.

\begin{example}[Reachability as a least fixed point]
Let $R\subseteq X\times X$ be a relation on a set and $S\subseteq X$.
Use graph image, $R[Q]=\{y\in X:\exists x\in Q,\ (x,y)\in R\}$.
This operation is monotone, so the preceding corollary gives the least
$R$-closed subset $C$ containing $S$, with $C=S\cup R[C]$.
It has the concrete description
\[
C=\bigcup_{n\in\mathbb N}C_n,\qquad
C_0=S,\quad C_{n+1}=C_n\cup R[C_n].
\]
Indeed, the union contains $S$. If $x$ belongs to it and $(x,y)\in R$,
then $x\in C_n$ for some $n$ and $y\in C_{n+1}$; hence the union is
$R$-closed. Conversely, induction shows that every $R$-closed set
containing $S$ contains every $C_n$, and therefore their union.
By induction on $n$, $C_n$ consists precisely of points reached from
$S$ by paths of at most $n$ edges: a new image point appends one edge,
and every nonzero-length path has a penultimate point. Thus $C$
records finite reachability, with zero-length paths retaining $S$.
This explicit countable iteration is justified for relational image;
monotonicity alone in an arbitrary complete lattice does not assert
that countably many iteration stages suffice.
\end{example}

%\ifdefined\ordertheoryflag
%\undef\ordertheoryflag
%\end{document}
%\fi
\ifdefined\qwertyflag
\undef\qwertyflag
\end{document}
\fi
